% --- Archon physics formalization source begin ---
% archon:physics
% archon:covers IPhO2026Problems/problem_IPhO_2026_1_C_1.lean
% archon:source-report reports/ipho_2026_k3/problem_IPhO_2026_1_C_1.source.json
% archon:problem-id IPhO_2026_1
% archon:part-id C.1

\chapter{Physics problem IPhO\_2026\_1\_C\_1}
\label{ch:IPhO2026Problems_problem_IPhO_2026_1_C_1}

\paragraph{Problem source.}
A photon of angular frequency omega is absorbed by an ozone molecule O3 at rest,
dissociating it into O2 and O.  Let U\_i and U\_f be the ground-state energies of
O3 and O2 and define Delta U = U\_f - U\_i.  The outgoing O2 momentum makes angle
theta with the incident photon.  Treat the oxygen fragments classically and
non-relativistically, take the mass of an oxygen atom to be m, and use photon
momentum p\_gamma = E\_gamma/c = hbar*omega/c.

Current subquestion:
Determine the minimum angular frequency omega\_min required for dissociation at outgoing O2 angle theta, in terms of hbar, c, theta, Delta U, and m.

\paragraph{Current subquestion.}
Determine the minimum angular frequency omega\_min required for dissociation at outgoing O2 angle theta, in terms of hbar, c, theta, Delta U, and m.

\paragraph{Recorded answer/context.}
For theta <= pi/2, omega\_min = 3*m*c\textasciicircum{}2*[1 - sqrt(1 - (Delta U/(3*m*c\textasciicircum{}2))*(2*sin(theta)\textasciicircum{}2 + 1))]/[hbar*(2*sin(theta)\textasciicircum{}2 + 1)]. For theta >= pi/2 use the same threshold evaluated at theta = pi/2.

\paragraph{Figure/image path.}
/root/proposal\_for\_physic/science-mango/ipho\_2026\_source/image/T1\_page-3.png

\paragraph{Formalization target.}
create a compiling Lean file with sorry bodies at `IPhO2026Problems/problem\_IPhO\_2026\_1\_C\_1.lean`.
The Lean declarations must preserve the physical quantities, dimensions or dimensional roles, figure labels, governing-law hypotheses, and final relation expressed by this problem.
Use Mathlib/Physlib names found through LeanExplore where available. If a domain API is missing, introduce faithful local abstractions rather than scalar placeholder aliases.

\begin{theorem}[Physics formalization target]
\label{thm:physics:IPhO_2026_1_C_1:target}
\uses{thm:IPhO2026Problems_problem_IPhO_2026_1_C_1:minimum_angular_frequency_T1_C1, thm:IPhO2026Problems_problem_IPhO_2026_1_C_1:hbarOmegaMin_pi_sub}
The assigned autoformalize agent should translate this physics problem into Lean declarations in the covered file, with theorem and lemma proof bodies written as `by sorry`.
\end{theorem}
\begin{proof}
This is an autoformalization task, not a proof task. Produce faithful statements that can later be proved without weakening the source contract.
\end{proof}
% NOTE: PhysLean-coverage exemption (planner-recorded, iter-002): PhysLean has no module for relativistic two-body photodissociation kinematics (see this file's physics-grounding log for the near-miss query results). Self-containment is kept with the `import Mathlib` baseline; no irrelevant Physlib import is added. This documents the resolved import policy (iter-001 review ruling) for the blueprint-doctor `missing-physlib-import` check.
% --- Archon physics formalization source end ---

% --- Archon named-quantities coverage (blueprint-writer 1-c-1-entries) ---
% NOTE: import policy (mirrors the reconciliation note above, iter-002 exemption):
% the covered file builds on the `import Mathlib` baseline alone; PhysLean has no
% module for relativistic two-body photodissociation kinematics, so no Physlib
% import is added.  This NOTE records the resolved import policy for the chapter.
\subsection*{Quantities and data}

\begin{definition}[Dissociation energy gap]
\label{def:IPhO2026Problems_problem_IPhO_2026_1_C_1:dissociationEnergyGap}
\lean{IPhO2026.Problem1.C1.dissociationEnergyGap}
The dissociation gap of the ozone breakup, $\Delta U = U_f - U_i$: the
difference of the ground-state energy of the fragments $O_2 + O$ and of
the ozone molecule $O_3$, which are opaque constants of the formalization.
\end{definition}
\begin{proof}
Definition; no claim.
\end{proof}

\begin{definition}[Positive-constant regime]
\label{def:IPhO2026Problems_problem_IPhO_2026_1_C_1:ConstantRegime}
\lean{IPhO2026.Problem1.C1.ConstantRegime}
\uses{def:IPhO2026Problems_problem_IPhO_2026_1_C_1:dissociationEnergyGap}
The physical regime of the problem: the reduced Planck constant $\hbar$,
the speed of light $c$, and the oxygen-atom mass $m$ are positive, and the
dissociation is endothermic, $\Delta U > 0$.
\end{definition}
\begin{proof}
Definition; no claim.
\end{proof}

\subsection*{Reaction geometry}

\begin{definition}[Reaction plane]
\label{def:IPhO2026Problems_problem_IPhO_2026_1_C_1:ReactionPlane}
\lean{IPhO2026.Problem1.C1.ReactionPlane}
The plane containing the incident photon and the outgoing fragments of
Figure~1c, modelled as an abstract two-dimensional real Euclidean space so
that the angle between momentum vectors is intrinsic.
\end{definition}
\begin{proof}
Definition; no claim.
\end{proof}

\begin{definition}[Incident photon line]
\label{def:IPhO2026Problems_problem_IPhO_2026_1_C_1:PhotonLine}
\lean{IPhO2026.Problem1.C1.PhotonLine}
\uses{def:IPhO2026Problems_problem_IPhO_2026_1_C_1:ReactionPlane}
The direction of the incident photon (and of its momentum) in the reaction
plane: a vector $\hat{k}$ certified to be a unit vector,
$\|\hat{k}\| = 1$.
\end{definition}
\begin{proof}
Definition; no claim.
\end{proof}

\begin{definition}[Scattering angle readout]
\label{def:IPhO2026Problems_problem_IPhO_2026_1_C_1:IsScatteringAngle}
\lean{IPhO2026.Problem1.C1.IsScatteringAngle}
\uses{def:IPhO2026Problems_problem_IPhO_2026_1_C_1:ReactionPlane, def:IPhO2026Problems_problem_IPhO_2026_1_C_1:PhotonLine}
The Figure-1c readout that $\theta$ is the angle between the outgoing
$O_2$ momentum $p$ and the incident photon direction $\hat{k}$: the
fragment is strictly outgoing, $p \neq 0$, and the cosine-law component
statement $\langle \hat{k}, p \rangle = \|p\| \cos\theta$ holds.
\end{definition}
\begin{proof}
Definition; no claim.
\end{proof}

\begin{definition}[Angular range of the problem]
\label{def:IPhO2026Problems_problem_IPhO_2026_1_C_1:IsAngularRange}
\lean{IPhO2026.Problem1.C1.IsAngularRange}
The figure range of the scattering angle: $\theta \in [0, \pi]$.  Branch
information is preserved by hypothesis, not selected in conclusions.
\end{definition}
\begin{proof}
Definition; no claim.
\end{proof}

\begin{definition}[Forward branch]
\label{def:IPhO2026Problems_problem_IPhO_2026_1_C_1:IsForwardBranch}
\lean{IPhO2026.Problem1.C1.IsForwardBranch}
The forward, non-backscattering branch of the official solution:
$\theta \le \pi/2$, the branch in which the tangent-line critical
configuration of Figure~1c exists.
\end{definition}
\begin{proof}
Definition; no claim.
\end{proof}

\subsection*{Governing-law structures}

\begin{definition}[Lawful two-body dissociation]
\label{def:IPhO2026Problems_problem_IPhO_2026_1_C_1:IsTwoBodyDissociation}
\lean{IPhO2026.Problem1.C1.IsTwoBodyDissociation}
\uses{def:IPhO2026Problems_problem_IPhO_2026_1_C_1:PhotonLine, def:IPhO2026Problems_problem_IPhO_2026_1_C_1:ReactionPlane, def:IPhO2026Problems_problem_IPhO_2026_1_C_1:IsScatteringAngle}
The governing-law structure of the absorption and two-body breakup
$\gamma + O_3 \to O_2 + O$ at photon angular frequency $\omega$: vector
momentum conservation fixes the oxygen-atom momentum
$q = (\hbar\omega/c)\,\hat{k} - p$; its eliminated cosine-law shadow is
$\|q\|^2 = (\hbar\omega/c)^2 + \|p\|^2 - 2(\hbar\omega/c)\,\|p\|\cos\theta$;
$\theta$ is the Figure-1c scattering angle; and non-relativistic energy
balance reads $\hbar\omega = \Delta U + \|p\|^2/(2\cdot 2m) + \|q\|^2/(2m)$
with fragment masses $2m$ and $m$.  No field mentions the threshold value.
\end{definition}
\begin{proof}
Definition; no claim.
\end{proof}

\begin{lemma}[Cosine-law elimination from vector momentum balance]
\label{lem:IPhO2026Problems_problem_IPhO_2026_1_C_1:momentum_q_sq_of_vector_balance}
\lean{IPhO2026.Problem1.C1.momentum_q_sq_of_vector_balance}
\uses{def:IPhO2026Problems_problem_IPhO_2026_1_C_1:PhotonLine, def:IPhO2026Problems_problem_IPhO_2026_1_C_1:IsScatteringAngle}
If $q = (\hbar\omega/c)\,\hat{k} - p$ and $\theta$ is the scattering angle
between $\hat{k}$ and $p$, then
$\|q\|^2 = (\hbar\omega/c)^2 + \|p\|^2 - 2(\hbar\omega/c)\,\|p\|\cos\theta$:
the scalar magnitude field of the dissociation structure is a genuine
consequence of the vector law, not an independent modelling assumption.
\end{lemma}
\begin{proof}
Expand $\|q\|^2 = \langle q, q \rangle$ with $q$ given by the vector
balance, use $\|\hat{k}\| = 1$ for the quadratic term, and insert the
readout $\langle \hat{k}, p \rangle = \|p\|\cos\theta$ for the cross term.
\end{proof}

\begin{definition}[Reachable frequency]
\label{def:IPhO2026Problems_problem_IPhO_2026_1_C_1:ReachableFrequency}
\lean{IPhO2026.Problem1.C1.ReachableFrequency}
\uses{def:IPhO2026Problems_problem_IPhO_2026_1_C_1:IsTwoBodyDissociation}
A photon angular frequency $\omega$ is reachable at outgoing $O_2$ angle
$\theta$ (atom mass $m$, dissociation gap $\Delta U$) iff some lawful
dissociation configuration realizes that frequency: an existential over
photon directions and fragment momenta, not a formula.
\end{definition}
\begin{proof}
Definition; no claim.
\end{proof}

\begin{definition}[Dissociation threshold]
\label{def:IPhO2026Problems_problem_IPhO_2026_1_C_1:IsDissociationThreshold}
\lean{IPhO2026.Problem1.C1.IsDissociationThreshold}
\uses{def:IPhO2026Problems_problem_IPhO_2026_1_C_1:ReachableFrequency}
$\omega_0$ is a dissociation threshold at angle $\theta$ iff it is
reachable and no strictly smaller positive frequency is reachable: a pure
minimality property that leaves the numerical value of the threshold open.
\end{definition}
\begin{proof}
Definition; no claim.
\end{proof}

\subsection*{Threshold candidate and derivation bridges}

\begin{definition}[Candidate threshold expression]
\label{def:IPhO2026Problems_problem_IPhO_2026_1_C_1:hbarOmegaMin}
\lean{IPhO2026.Problem1.C1.hbarOmegaMin}
The candidate closed-form forward-branch threshold
$\Omega(m, c, \Delta U, \theta)$ of the official solution, named as a bare
real-valued expression of the problem data (built from the angular factor
$2\sin^2\theta + 1$ and the square root of its discriminant).  By itself it
asserts nothing about dissociation; the claim that it \emph{is} the minimum
reachable frequency is the content of the target theorems below.
\end{definition}
\begin{proof}
Definition; no claim.
\end{proof}

\begin{lemma}[Angular-factor identity]
\label{lem:IPhO2026Problems_problem_IPhO_2026_1_C_1:two_sin_sq_add_one_eq}
\lean{IPhO2026.Problem1.C1.two_sin_sq_add_one_eq}
For every angle $\theta$, $2\sin^2\theta + 1 = 2 - \cos 2\theta$: the
angular factor of the threshold formula equals the coefficient appearing in
the official quadratic characterization.
\end{lemma}
\begin{proof}
The double-angle identity $\cos 2\theta = 1 - 2\sin^2\theta$ rearranged.
\end{proof}

\begin{lemma}[Quadratic characterization of the candidate threshold]
\label{lem:IPhO2026Problems_problem_IPhO_2026_1_C_1:quadratic_characterization_of_threshold}
\lean{IPhO2026.Problem1.C1.quadratic_characterization_of_threshold}
\uses{def:IPhO2026Problems_problem_IPhO_2026_1_C_1:hbarOmegaMin, def:IPhO2026Problems_problem_IPhO_2026_1_C_1:IsForwardBranch, lem:IPhO2026Problems_problem_IPhO_2026_1_C_1:two_sin_sq_add_one_eq}
For positive $m, c, \Delta U$, a forward angle $0 < \theta \le \pi/2$ with
nondegenerate angular factor $2 - \cos 2\theta \neq 0$ and a real square
root in the candidate formula, the scaled candidate
$E_0 = \hbar\,\Omega(m,c,\Delta U,\theta)$ is positive, solves
$(2 - \cos 2\theta)\,E^2 - 6 m c^2 E + 6\,\Delta U\,m c^2 = 0$,
and is smaller than every other positive root: it is the smallest positive
root of the official-solution quadratic.
\end{lemma}
\begin{proof}
Rewrite $2\sin^2\theta + 1$ as $2 - \cos 2\theta$ via the angular-factor
identity, so the candidate becomes the root of the displayed quadratic in
$E = \hbar\omega$ selected by the minus sign of the quadratic formula;
positivity follows from the positive regime and the real square root, and
minimality among positive roots because the two roots are ordered by that
sign choice.
\end{proof}

\subsection*{Value and target theorems}

\begin{theorem}[Minimum angular frequency, forward branch]
\label{thm:IPhO2026Problems_problem_IPhO_2026_1_C_1:minimum_angular_frequency_T1_C1}
\lean{IPhO2026.Problem1.C1.minimum_angular_frequency_T1_C1}
\uses{def:IPhO2026Problems_problem_IPhO_2026_1_C_1:ConstantRegime, def:IPhO2026Problems_problem_IPhO_2026_1_C_1:IsAngularRange, def:IPhO2026Problems_problem_IPhO_2026_1_C_1:IsForwardBranch, def:IPhO2026Problems_problem_IPhO_2026_1_C_1:IsDissociationThreshold, def:IPhO2026Problems_problem_IPhO_2026_1_C_1:hbarOmegaMin, lem:IPhO2026Problems_problem_IPhO_2026_1_C_1:quadratic_characterization_of_threshold, lem:IPhO2026Problems_problem_IPhO_2026_1_C_1:momentum_q_sq_of_vector_balance}
In the positive-constant regime, for a nondegenerate forward angle
$0 < \theta \le \pi/2$ with a real square root, the candidate expression
\[
\Omega(m,c,\Delta U,\theta)
= \frac{3 m c^2 \left(1 - \sqrt{1 - \frac{\Delta U}{3 m c^2}\,(2\sin^2\theta + 1)}\right)}{\hbar\,(2\sin^2\theta + 1)}
\]
is the dissociation threshold at $\theta$: some lawful two-body
configuration realizes it, and no smaller positive frequency is reachable.
The recorded closed form appears strictly on the conclusion side; the
hypotheses carry only the physical regime, the figure angle range, and the
algebraic nondegeneracy.
\end{theorem}
\begin{proof}
Eliminate $\|p\|$ between the energy-balance and cosine-law fields of a
lawful configuration at frequency $\omega$: reachability at forward angle
$\theta$ reduces to the official-solution quadratic in $\hbar\omega$
having a real solution, which holds exactly at and above its smallest
positive root.  By the quadratic characterization that root is
$\hbar\,\Omega$, so $\Omega$ is reachable and minimal.
\end{proof}

\begin{theorem}[Minimum angular frequency, backward branch]
\label{thm:IPhO2026Problems_problem_IPhO_2026_1_C_1:minimum_angular_frequency_backward_branch_T1_C1}
\lean{IPhO2026.Problem1.C1.minimum_angular_frequency_backward_branch_T1_C1}
\uses{def:IPhO2026Problems_problem_IPhO_2026_1_C_1:ConstantRegime, def:IPhO2026Problems_problem_IPhO_2026_1_C_1:IsAngularRange, def:IPhO2026Problems_problem_IPhO_2026_1_C_1:IsDissociationThreshold, def:IPhO2026Problems_problem_IPhO_2026_1_C_1:hbarOmegaMin, thm:IPhO2026Problems_problem_IPhO_2026_1_C_1:minimum_angular_frequency_T1_C1, thm:IPhO2026Problems_problem_IPhO_2026_1_C_1:hbarOmegaMin_pi_sub}
In the positive-constant regime, for a backward angle
$\pi/2 \le \theta \le \pi$ with a real square root (at $\theta = \pi/2$),
the dissociation threshold freezes at its $\theta = \pi/2$ value: the
threshold at $\theta$ is the forward-branch candidate evaluated at
$\pi/2$, $\Omega(m, c, \Delta U, \pi/2)$.  The recorded answer again
occurs only in the conclusion.
\end{theorem}
\begin{proof}
The tangent-line critical configuration of Figure~1c exists only in the
forward branch, so for $\theta \ge \pi/2$ the binding constraint is the
$\pi/2$ one; formally this rests on the forward-branch threshold theorem
at $\pi/2$ together with the reflection symmetry of the candidate formula.
\end{proof}

\begin{theorem}[Reflection symmetry of the candidate threshold]
\label{thm:IPhO2026Problems_problem_IPhO_2026_1_C_1:hbarOmegaMin_pi_sub}
\lean{IPhO2026.Problem1.C1.hbarOmegaMin_pi_sub}
\uses{def:IPhO2026Problems_problem_IPhO_2026_1_C_1:hbarOmegaMin}
The candidate threshold is invariant under the reflection
$\theta \mapsto \pi - \theta$:
$\Omega(m, c, \Delta U, \pi - \theta) = \Omega(m, c, \Delta U, \theta)$,
the pure-math symmetry justifying the backward-branch freeze.
\end{theorem}
\begin{proof}
The candidate expression depends on $\theta$ only through $\sin^2\theta$,
and $\sin(\pi - \theta) = \sin\theta$.
\end{proof}
